\chapter{Contexto, Diagnóstico del Sistema Heredado y Objetivos}

\section{Información del Cliente Comercial y Problemática}
El presente proyecto de ingeniería de software responde a las necesidades operativas y comerciales de la tienda de indumentaria masculina \textbf{Machi's}, ubicada en la ciudad de Tacna, Perú. La empresa cuenta actualmente con dos establecimientos comerciales físicos en zonas de alto tránsito de la ciudad:
\begin{enumerate}
    \item \textbf{Local Principal:} Ubicado en la Av. San Martín N$^\circ$ 450.
    \item \textbf{Sucursal Mall:} Ubicada en la Av. Bolognesi N$^\circ$ 820.
\end{enumerate}

\subsection{Diagnóstico del Negocio}
A pesar de su posicionamiento en el mercado local de ropa masculina, el negocio operaba con serias limitaciones de gestión:
\begin{itemize}
    \item \textbf{Desconexión de inventarios:} No existía un canal centralizado en tiempo real para consultar existencias entre ambas tiendas, provocando pérdidas de ventas cuando un cliente solicitaba una talla disponible en la otra sucursal.
    \item \textbf{Inconsistencia en registros de prendas:} El uso de registros manuales y planillas de cálculo sin listas cerradas permitía el ingreso arbitrario de tallas (ej. ``M'', ``m'', ``Mediano'') y colores (ej. ``Azul'', ``Marino'', ``Blue''), corrompiendo la base de datos y distorsionando los balances de existencias.
    \item \textbf{Ausencia de presencia digital:} Los clientes finales carecían de un catálogo web para consultar prendas, precios y disponibilidad en tienda antes de acercarse al establecimiento físico.
    \item \textbf{Emisión desarticulada de comprobantes:} No existía emisión digital de tickets de control interno, dificultando el arqueo de caja al cierre de jornada.
\end{itemize}

\section{Diagnóstico Crítico del Sistema Heredado (Análisis de Sistemas v1.0)}
En el curso previo de \textit{Análisis de Sistemas} (documentado en el informe \textit{Trabajo Final II}), se realizó un primer acercamiento empírico al proyecto. La evaluación crítica de dicho antecedente reveló deficiencias estructurales que hacían inviable su paso directo a producción:

\begin{enumerate}
    \item \textbf{Atomización extrema y confusión de niveles de abstracción:} En el informe previo, cada caja de texto o campo de formulario de interfaz gráfica fue tipificado erróneamente como un ``Requisito Funcional'' independiente (ej. \textit{RF-08: Pedir nombre de usuario}, \textit{RF-09: Pedir contraseña}, \textit{RF-10: Color de prenda}, \textit{RF-11: Talla de prenda}). Esto infló artificialmente un catálogo de 72 requerimientos superficiales carentes de lógica transaccional, reglas de negocio y restricciones de integridad.
    \item \textbf{Vulnerabilidad crítica por ausencia de catálogos maestros cerrados:} El sistema permitía la escritura libre de atributos clave. En una tienda de ropa comercial, la introducción de texto abierto en tallas y colores imposibilita el control de stock por SKU y quiebra los filtros de búsqueda.
    \item \textbf{Confusión regulatoria y riesgo fiscal (SUNAT):} El análisis heredado hacía mención de ``facturación y boletas electrónicas'', generando la falsa presunción de integración con servidores tributarios (SUNAT/OSE), lo cual resultaba sobredimensionado y legalmente riesgoso para un negocio minorista que requería exclusivamente \textbf{tickets digitales de control interno}.
    \item \textbf{Omisión de procesos transaccionales indispensables:} El sistema previo omitió el control de turnos de caja (apertura, arqueo ciego y cierre), el soporte de pagos mixtos (Efectivo + billeteras digitales Yape/Plin), el control de mermas en Kardex y la anulación supervisada de ventas.
    \item \textbf{Criterios de aceptación no verificables:} Criterios como \textit{``Se acepta cuando funcione sin errores''} violaban las buenas prácticas de calidad de software y testing formal.
    \item \textbf{Monopolio de roles y error en la EDT (WBS):} Los cinco paquetes de trabajo de programación estaban asignados a un único responsable con una codificación errónea (\texttt{1.3.1.X} en lugar de \texttt{1.3.3.X}), sin contemplar especialización de roles ni arquitectura de software.
\end{enumerate}

\begin{figure}[H]
\centering
\begin{tikzpicture}[
    scale=0.92, every node/.style={transform shape},
    boxLegacy/.style={rectangle, draw=red!80!black, fill=red!10, rounded corners=3pt, thick, align=center, minimum width=6cm, minimum height=0.7cm},
    boxModern/.style={rectangle, draw=blue!80!black, fill=blue!10, rounded corners=3pt, thick, align=center, minimum width=6cm, minimum height=0.7cm},
    headTitle/.style={font=\bfseries\normalsize, align=center}
]
% Lado Legacy
\node[headTitle, red!85!black] (tLegacy) at (0, 3.8) {Sistema Heredado (Legacy v1.0)\\\small Enfoque Empírico / Acoplado};
\node[boxLegacy] (l1) at (0, 2.7) {UI Monolítica (Vistas acopladas a DB)};
\node[boxLegacy] (l2) at (0, 1.8) {Lógica Dispersa (Sin capas ni CQRS)};
\node[boxLegacy] (l3) at (0, 0.9) {DB 5 Tablas (Desnormalizada, campos abiertos)};
\node[boxLegacy] (l4) at (0, 0.0) {Testing Subjetivo (Sin casos formales)};

% Flecha de Transformacion
\draw[-{Stealth[length=3.5mm]}, line width=2pt, primaryblue] (3.4, 1.35) -- (5.0, 1.35) 
    node[midway, above=3pt, font=\bfseries\small, text=primaryblue] {Refactorización}
    node[midway, below=3pt, font=\scriptsize\itshape, text=darkblue] {Ingeniería v2.0};

% Lado Moderno
\node[headTitle, blue!85!black] (tModern) at (8.4, 3.8) {Solución Profesional (SYSELL v2.0)\\\small Clean Architecture + Estándares};
\node[boxModern, fill=green!15, draw=green!70!black] (m1) at (8.4, 2.7) {Capa Presentación (React 18 + SPA Mobile)};
\node[boxModern, fill=cyan!15, draw=cyan!70!black] (m2) at (8.4, 1.8) {Capa Aplicación \& Dominio (CQRS + RBAC)};
\node[boxModern, fill=orange!15, draw=orange!70!black] (m3) at (8.4, 0.9) {MySQL 8.0 en 3FN (17 Tablas normalizadas)};
\node[boxModern, fill=purple!15, draw=purple!70!black] (m4) at (8.4, 0.0) {SQA IEEE 829 (77 CPs) + CI/CD Docker};

\end{tikzpicture}
\caption{Comparativa Estructural: De Enfoque Empírico Legacy a Ingeniería de Software v2.0}
\label{fig:legacy_vs_modern}
\end{figure}

\section{Matriz de Transformación Integral: Análisis de Sistemas $\rightarrow$ v2.0}
Para superar las debilidades del proyecto heredado y elevar la solución a los estándares profesionales de ingeniería de software, se ejecutó una refundación metodológica integral estructurada en once dimensiones técnicas:

\begin{table}[H]
\centering
\small
\begin{tabularx}{\textwidth}{>{\bfseries}p{2.6cm} Y Y >{\centering\arraybackslash}p{2.3cm}}
\toprule
\textbf{Capa de Ingeniería} & \textbf{Análisis de Sistemas (v1.0)} & \textbf{Ingeniería de Software (v2.0)} & \textbf{Criterio Rúbrica} \\
\midrule
Requisitos & 72 requisitos empíricos atomizados por campos de pantalla. & \textbf{77 Requisitos Funcionales Oficiales} (ISO 29148) estructurados con criterios verificables. & Crit. 6, 7, 12 \\
\addlinespace
Calidad / RNF & Requisitos no funcionales genéricos sin métricas. & \textbf{14 Requisitos No Funcionales (RNF)} categorizados bajo norma \textbf{ISO/IEC 25010}. & Crit. 6, 15 \\
\addlinespace
Trazabilidad & Inexistente; requisitos huérfanos sin correspondencia. & \textbf{Matriz Integral de 7 columnas} con cobertura bidireccional al 100\%. & Crit. 8, 14, 16 \\
\addlinespace
Procesos (DFD) & DFDs conceptuales desconectados de almacenes. & \textbf{DFD Niveles 0, 1 y 2 balanceados} con almacenes normalizados \texttt{D1} a \texttt{D5}. & Crit. 11 \\
\addlinespace
Casos de Uso & Inexistente (cero diagramas formales UML). & \textbf{15 Casos de Uso UML 2.5} con especificaciones textuales y flujos alternos. & Crit. 11, 12 \\
\addlinespace
Diseño Dinámico & Cero diagramas de interacción. & \textbf{8 Diagramas de Secuencia PlantUML} (\texttt{strictuml}) con transacciones ACID. & Crit. 11, 16 \\
\addlinespace
Arquitectura & Estructura monolítica no documentada. & \textbf{Clean Architecture (4 capas) + CQRS} y autenticación híbrida Laravel Sanctum. & Crit. 2, 3, 4 \\
\addlinespace
Base de Datos & Diseño empírico informal de tablas aisladas. & \textbf{17 Tablas relacionales en 3FN} con scripts \texttt{schema.sql} y \texttt{seeds.sql}. & Crit. 11, 18 \\
\addlinespace
Testing / QA & ``Se acepta si funciona sin errores'' (subjetivo). & \textbf{77 Casos de Prueba IEEE 829 / ISTQB} con validación positiva, negativa y límites. & Crit. 10, 15 \\
\addlinespace
Gestión Ágil & EDT en cascada rígida con errores jerárquicos. & \textbf{Híbrido RUP + Scrum:} Product Backlog (18 HUs, 94 SP) en 4 Sprints con DoR/DoD. & Crit. 4, 5, 9, 10 \\
\addlinespace
DevOps / CI & Despliegue manual no documentado. & \textbf{Docker Compose} y \textbf{CI en GitHub Actions} con validación continua. & Crit. 18, 20 \\
\bottomrule
\end{tabularx}
\caption{Matriz comparativa de evolución capa por capa: Análisis de Sistemas vs. v2.0}
\label{tab:matriz_transformacion}
\end{table}

\section{Objetivos del Proyecto}

\subsection{Objetivo General}
Diseñar, modelar, especificar e implementar un sistema web integral de gestión comercial, control de inventario multilocal, terminal punto de venta (POS) y vitrina digital pública responsiva (\textbf{SYSELL Web v2.0}) para la tienda comercial de indumentaria masculina \textit{Machi's}, satisfaciendo los más rigurosos estándares de arquitectura limpia, trazabilidad y aseguramiento de la calidad de software.

\subsection{Objetivos Específicos}
\begin{enumerate}
    \item Implementar la gestión de existencias centralizada con control de inventario independiente para los dos locales comerciales, alertas de stock mínimo y registro auditable de movimientos en Kardex.
    \item Desarrollar una terminal de punto de venta rápida (POS) con soporte para pagos multitendencia mixtos (Efectivo, tarjetas y billeteras digitales Yape/Plin), cálculo de cambio y emisión de tickets digitales con trazabilidad interna desligada de SUNAT.
    \item Construir una plataforma de vitrina virtual pública responsiva que permita a los clientes finales consultar en tiempo real las prendas en stock, precios, fotografías en alta definición y disponibilidad física discriminada por sucursal.
    \item Garantizar la integridad y consistencia del software mediante una arquitectura desacoplada en 4 capas concéntricas (Clean Architecture) con patrón CQRS, respaldada por una cobertura del 100\% de pruebas funcionales basadas en el estándar IEEE 829.
\end{enumerate}
